\section{Discussion}
\label{sec:discussion}
%- Verdict, In distribution: The practical takeaway here is that for monitoring a judge, reading its internal states is more reliable than asking it or reading its confidence.
% Verdict shift: They are consequently more useful for prioritizing problems for verification than for selecting among candidates of the same problem.
%triage instead of selection since they help prioritize which problems to verify rather than select among a problem's candidates.

\paragraph{Verdict and program are different prediction targets}
Our results show that the usefulness of execution-free correctness prediction depends on both the prediction target and the type of distribution shift. In particular, predicting whether a model's verdict is correct is different from predicting whether the program itself is correct, even though both targets use signals from the same judging pass.
The two targets respond differently to shift. Causal localization helps only on the verdict, and between benchmarks only program correctness transfers. An execution-free correctness predictor for an LLM-based code-generation system should therefore be chosen and evaluated for the decision it supports. To decide which candidates to test in a new programming language, output confidence is a strong baseline: it needs no labels and loses almost nothing under shift,
while trained predictors keep only part of their lead. To decide whether to
trust the model's verdict, confidence is weak, and no verdict predictor we tested
carries over to a new benchmark. % without data from it.


% \paragraph{What a practitioner can take away}
% To decide whether an LLM's verdict on code can be trusted in a new
% programming language, a probe on a causally localized subspace was the most
% robust predictor we found, and once the subspace is learned it costs no more
% than any probe. The gain is model-dependent, however, and DAS's own quality
% measure does not reveal it: on gpt-oss, well-converged subspaces were no
% better than random ones. A practitioner therefore has to validate the subspace
% on held-out data from the target setting rather than trust IIA. Between
% benchmarks, no verdict predictor transferred reliably. The judge's errors
% change character: Qwen3-Coder accepts 58\% of the LCB candidates, of which
% 55\% pass, but 72\% of the BCB candidates, of which only 40\% pass
% (Table~\ref{tab:datasets}), so a predictor of its reliability learned on one
% benchmark meets a different error profile on the other. A few labeled
% examples from the target setting, which help probes in other domains
% considerably~\cite{McKenzie2025HighStakes}, are likely needed there.
% For program correctness, output confidence is a strong baseline that needs no
% labels and loses almost nothing under language shift; trained probes pay off
% in distribution, and much less after a shift.

\paragraph{Triage rather than selection}
Much of what the methods predict is which problems
the model solves: every method is weak at ranking candidates within a problem, and a probe that has seen only the problem matches one that has seen the candidate (Table \ref{tab:genprobe}). Pooled metrics,  such as AUROC~\cite{bouchard2026functional, shi2026codemorethantext} or accuracy and F1~\cite{bui2025openia, Vu2026Autoprobe}, thus mix problem difficulty with the ability to tell a correct candidate from an incorrect one. 
For triage across problems this is useful, while for choosing among
the candidates of one problem it is not. Selection metrics~\cite{ribeiro2026repe} and pairwise
accuracy within a task~\cite{ribeiro2026robustness} measure the second ability
directly; we recommend reporting such a within-problem measure alongside the
pooled one. 

\paragraph{Causal localization for code}
Huang et al.~\cite{Huang2025Internal} report that causally localized
predictors are more robust under shift on natural-language tasks. For code,
this holds on one of our two models and for the target the method was
designed for, the model's own verdict; on the program correctness target it gives no
advantage. Where it helps, the learned subspace carries the gain: once the subspace is
learned, a probe on it is as robust as simulating interventions and costs no
more than any probe. IIA does not reliably identify subspaces that are useful for correctness prediction. On gpt-oss, well-converged subspaces were no better than random ones, so the
subspace has to be validated on held-out data.
%, not the intervention: reading the subspace with a probe is as robust as simulating interventions, at a fraction of the cost. %\orange{Why the subspace helps on Qwen3-Coder and not on gpt-oss we can only partly explain. On gpt-oss, the useful subspaces are those at layers where DAS has not converged; one possibility is that gpt-oss, a reasoning model judging without reasoning, has fully formed its decision at these layers, so that the subspace holds the decision but little about whether it is right.}